\textbf{Sampling and the Nyquist Frequency}

\vspace{0.5cm}

% Continuous signal with adequate sampling
\begin{tikzpicture}
\begin{axis}[
    width=13cm, height=4cm,
    xlabel={Time $t$ (seconds)},
    ylabel={Amplitude},
    xmin=0, xmax=2,
    ymin=-1.3, ymax=1.3,
    grid=major,
    grid style={gray!30},
    title={\normalsize Adequate Sampling: $f_s = 10$ Hz $> 2f_{\text{signal}} = 6$ Hz (Nyquist criterion satisfied)},
    title style={at={(0.5,1.05)}, anchor=south}
]

% Original continuous signal: 3 Hz sine wave
\addplot[UofSC50Black, dashed, thick, domain=0:2, samples=500]
    {sin(deg(2*pi*3*x))};

% Sample points at 10 Hz (every 0.1 s)
\addplot+[UofSCAtlantic, only marks, mark=*, mark size=3pt]
    coordinates {
        (0, 0)
        (0.1, 0.951)
        (0.2, 0.588)
        (0.3, -0.588)
        (0.4, -0.951)
        (0.5, 0)
        (0.6, 0.951)
        (0.7, 0.588)
        (0.8, -0.588)
        (0.9, -0.951)
        (1.0, 0)
        (1.1, 0.951)
        (1.2, 0.588)
        (1.3, -0.588)
        (1.4, -0.951)
        (1.5, 0)
        (1.6, 0.951)
        (1.7, 0.588)
        (1.8, -0.588)
        (1.9, -0.951)
        (2.0, 0)
    };

% Reconstructed signal (matches original well)
\addplot[UofSCAtlantic, very thick, domain=0:2, samples=500]
    {sin(deg(2*pi*3*x))};

% Legend entries
\addlegendentry{Original: $f = 3$ Hz}
\addlegendentry{Sample points}
\addlegendentry{Reconstructed}

% Annotation
\node[UofSCHorseshoe, font=\footnotesize, fill=white, inner sep=2pt] at (1.5, 0.85) {$f_s > 2f_{\text{max}}$ : No aliasing};

\end{axis}
\end{tikzpicture}

\vspace{0.6cm}

% Undersampled signal showing aliasing
\begin{tikzpicture}
\begin{axis}[
    width=13cm, height=4cm,
    xlabel={Time $t$ (seconds)},
    ylabel={Amplitude},
    xmin=0, xmax=2,
    ymin=-1.3, ymax=1.3,
    grid=major,
    grid style={gray!30},
    title={\normalsize Undersampling: $f_s = 4$ Hz $< 2f_{\text{signal}} = 6$ Hz (Aliasing occurs!)},
    title style={at={(0.5,1.05)}, anchor=south}
]

% Original continuous signal: 3 Hz sine wave
\addplot[UofSC50Black, dashed, thick, domain=0:2, samples=500]
    {sin(deg(2*pi*3*x))};

% Sample points at 4 Hz (every 0.25 s) - sampling a 3 Hz signal
% sin(2*pi*3*t) at t = 0, 0.25, 0.5, 0.75, 1.0, ...
\addplot+[UofSCGarnet, only marks, mark=*, mark size=3pt]
    coordinates {
        (0, 0)
        (0.25, -0.707)
        (0.5, 0)
        (0.75, 0.707)
        (1.0, 0)
        (1.25, -0.707)
        (1.5, 0)
        (1.75, 0.707)
        (2.0, 0)
    };

% Aliased signal: appears as 1 Hz signal (|3 - 4| = 1 Hz alias)
\addplot[UofSCRose, very thick, domain=0:2, samples=300]
    {-sin(deg(2*pi*1*x))};

% Legend
\addlegendentry{Original: $f = 3$ Hz}
\addlegendentry{Sample points}
\addlegendentry{Aliased: $f_{\text{alias}} = |f - f_s| = 1$ Hz}

% Annotation
\node[UofSCGarnet, font=\footnotesize, fill=white, inner sep=2pt] at (1.5, -0.85) {$f_s < 2f_{\text{max}}$ : \textbf{Aliasing!}};

\end{axis}
\end{tikzpicture}

\vspace{0.6cm}

% Frequency domain illustration
\begin{tikzpicture}
\begin{axis}[
    width=6.2cm, height=4.5cm,
    xlabel={Frequency $f$ (Hz)},
    ylabel={$|X(f)|$},
    xmin=-8, xmax=8,
    ymin=0, ymax=1.3,
    grid=major,
    grid style={gray!30},
    title={\normalsize Adequate Sampling Spectrum},
    title style={at={(0.5,1.05)}, anchor=south},
    xtick={-6,-4,-3,-2,0,2,3,4,6},
    ytick={0, 0.5, 1}
]

% Original spectrum (impulses at +/- 3 Hz)
\addplot+[UofSCAtlantic, very thick, ycomb, mark=triangle*, mark options={fill=UofSCAtlantic, scale=1.5}]
    coordinates {(-3, 1) (3, 1)};

% Nyquist frequency markers
\draw[UofSCHorseshoe, dashed, very thick] (-5, 0) -- (-5, 1.2);
\draw[UofSCHorseshoe, dashed, very thick] (5, 0) -- (5, 1.2);

% Labels
\node[UofSCAtlantic, font=\footnotesize, anchor=south] at (3, 1.02) {$f_0$};
\node[UofSCAtlantic, font=\footnotesize, anchor=south] at (-3, 1.02) {$-f_0$};
\node[UofSCHorseshoe, font=\footnotesize, anchor=south] at (5, 1.22) {$f_s/2$};
\node[UofSCHorseshoe, font=\footnotesize, anchor=south] at (-5, 1.22) {$-f_s/2$};

% Shade the Nyquist band
\fill[UofSCHorseshoe, opacity=0.1] (-5, 0) rectangle (5, 1.2);

\end{axis}
\end{tikzpicture}
\hfill
\begin{tikzpicture}
\begin{axis}[
    width=6.2cm, height=4.5cm,
    xlabel={Frequency $f$ (Hz)},
    ylabel={$|X(f)|$},
    xmin=-8, xmax=8,
    ymin=0, ymax=1.3,
    grid=major,
    grid style={gray!30},
    title={\normalsize Undersampled Spectrum (Aliasing)},
    title style={at={(0.5,1.05)}, anchor=south},
    xtick={-6,-4,-3,-2,-1,0,1,2,3,4,6},
    ytick={0, 0.5, 1}
]

% Original spectrum (impulses at +/- 3 Hz) - shown faded
\addplot+[UofSC50Black, thick, ycomb, mark=triangle*, mark options={fill=UofSC50Black, scale=1.2}]
    coordinates {(-3, 0.7) (3, 0.7)};

% Aliased frequencies (folded into Nyquist band)
% With fs=4, f_N=2, signal at 3 Hz aliases to |3-4|=1 Hz
\addplot+[UofSCRose, very thick, ycomb, mark=*, mark options={fill=UofSCRose, scale=1.5}]
    coordinates {(-1, 1) (1, 1)};

% Nyquist frequency markers
\draw[UofSCGarnet, dashed, very thick] (-2, 0) -- (-2, 1.2);
\draw[UofSCGarnet, dashed, very thick] (2, 0) -- (2, 1.2);

% Labels
\node[UofSC50Black, font=\footnotesize, anchor=south] at (3, 0.72) {$f_0$};
\node[UofSCRose, font=\footnotesize, anchor=south] at (1, 1.02) {$f_{\text{alias}}$};
\node[UofSCGarnet, font=\footnotesize, anchor=south] at (2, 1.22) {$f_s/2$};
\node[UofSCGarnet, font=\footnotesize, anchor=south] at (-2, 1.22) {$-f_s/2$};

% Shade the Nyquist band
\fill[UofSCGarnet, opacity=0.1] (-2, 0) rectangle (2, 1.2);

% Arrow showing folding
\draw[UofSCRose, thick, ->, >=stealth] (3, 0.5) to[out=180, in=0] (1, 0.7);
\node[UofSCRose, font=\tiny] at (2.3, 0.3) {Folds};

\end{axis}
\end{tikzpicture}

\vspace{0.4cm}

% Nyquist theorem statement box
\begin{tikzpicture}
\node[draw=UofSCGarnet, thick, fill=UofSCGarnet!8, minimum width=12.5cm, minimum height=1.8cm, sharp corners, inner sep=8pt, align=center] {
    \textbf{Nyquist-Shannon Sampling Theorem:} A bandlimited signal with maximum frequency $f_{\text{max}}$ \\[2pt]
    can be perfectly reconstructed from samples if $f_s > 2f_{\text{max}}$. \\[4pt]
    \textcolor{UofSCGarnet}{\textbf{Nyquist Frequency:}} $f_N = f_s/2$ \quad \textcolor{UofSCGarnet}{\textbf{Nyquist Rate:}} $f_s = 2f_{\text{max}}$
};
\end{tikzpicture}

